% Generated by report/make_restated.py from theory/A1_lower_bound.tex: do not edit (edit the note and rerun).
\begin{restated}{Theorem}{thm:LZ}{A1}{0}{Lan--Zhou lower bound; Theorem~3 and Corollary~3 of \cite{lan2018optimal}}
Let $\mathcal Q$ be the parameter of the instance \eqref{eq:LZinstance} and $q=\frac{\sqrt{\mathcal Q}-1}{\sqrt{\mathcal Q}+1}$. The iterates of any randomized incremental gradient method of class (S1) (Definition~\ref{def:span}) applied to \eqref{eq:LZinstance} satisfy
\[
\frac{\E\norm{x^k-x^\ast}^2}{\norm{x^0-x^\ast}^2}\ \ge\ \frac12\exp\Big(-\frac{4k\sqrt{\mathcal Q}}{m(\sqrt{\mathcal Q}+1)^2-4\sqrt{\mathcal Q}}\Big)
\qquad\text{for all }\ n\ge\underline n(m,k),
\]
\[
\underline n(m,k):=\frac{m\log\big[(1-(1-q^2)/m)^k/2\big]}{2\log q},
\]
where $n=m\tilde n$ is the total dimension. Consequently the number of gradient evaluations needed by such a method to find $\bar x$ with $\E\norm{\bar x-x^\ast}^2\le\varepsilon$ cannot be smaller than
\[
\Omega\Big\{\big(m+\sqrt{m\,\mathcal C}\big)\log\frac{\norm{x^0-x^\ast}^2}{\varepsilon}\Big\},\qquad \mathcal C=\frac{L}{\mu},\ L=\sum_{i=1}^mL_i ,
\]
if $n$ is sufficiently large; precisely (their display (3.45) in arXiv:1507.02000v3) $k\ge\big[\frac m2\big(\frac{\sqrt{\mathcal Q}}2+1\big)-1\big]\log\frac{\norm{x^0-x^\ast}^2}{2\varepsilon}$.
\end{restated}
